% Future work section for OrbitGraph paper.

\section{Future Work}
\label{sec:future}

Our study demonstrates that graph-centric SDN can eliminate GSL handover black
holes on emulated constellations up to 102 nodes, but several limitations
point to concrete extensions.
We group them by experimental methodology, scenario coverage, and path toward
production-scale deployment.

\subsection{Strengthening the evaluation harness}

\paragraph{Faster southbound interfaces.}
The largest gap between OrbitGraph's \emph{algorithmic} cost and its
\emph{observed} control-plane time is \texttt{install\_ms}: pushing routes
through per-container \texttt{docker exec} dominates handover and damage events
(Section~\ref{sec:analysis-control}).
Replacing this harness with a production-style southbound---gRPC or NETCONF to
on-board routing agents, compressed delta messages, or batched OpenFlow-style
flow mods---would isolate the true cost of centralized programming from Docker
overhead and likely improve damage-recovery performance
(Section~\ref{sec:analysis-damage}), where install latency currently outweighs
the benefit of global recomputation.

\paragraph{Separating install from the emulation clock.}
Our synchronous StarryNet loop blocks while routes are installed, which
inflates outage measurements for reactive events even when forwarding is
partially preserved.
Future work should run the controller asynchronously (install in background
threads or on a dedicated orchestration host) and align probes to wall-clock
control-plane timestamps only, yielding a fairer comparison with distributed
protocols whose reconvergence also proceeds concurrently with forwarding
attempts.

\paragraph{Richer data-plane workloads.}
The present campaign uses ICMP probes between two ground stations.
Adding sustained TCP/QUIC flows (\texttt{iperf3}), multipath traffic matrices,
and application-level QoS metrics would test whether zero handover outage
translates into throughput and goodput gains under realistic load, not only
reachability.

\paragraph{Extended statistical campaigns.}
Ten repetitions per grid size stabilize mean and standard deviation for our
scaling study, but seed sensitivity (especially for random satellite damage)
warrants wider confidence intervals and, where feasible, non-parametric tests
across the full $5{\times}5$–$10{\times}10$ matrix.
Automating regression checks on artifact completeness (missing handover phases,
truncated runs) would harden batch pipelines for larger parameter sweeps.

\subsection{Scenarios not yet covered}

\paragraph{Multiple handovers and long horizons.}
Each run captures the first GSL handover by design.
Mega-constellations experience continuous contact churn; emulating multi-hour
windows with dozens of handovers per ground station would reveal whether
incremental baselines stay consistent and whether error accumulation in address
caches or route-key bookkeeping appears over time.

\paragraph{Inter-domain and ground-relay paths.}
We use a single autonomous system with two fixed ground stations and
LeastDelay intra-domain routing.
Extending to multi-AS topologies with BGP peering at gateways, or to
ground-relay architectures that offload long-haul traffic to terrestrial
fiber~\cite{handley:ground-relays:2019}, would test OrbitGraph where path
selection crosses administrative boundaries and not every edge weight comes
from propagation delay alone.

\paragraph{Comparison with orbit-aware distributed routing.}
OrbitGraph proactive handover is complementary to, not mutually exclusive with,
distributed protocols that exploit ephemeris.
A head-to-head study against OPSPF~\cite{pam:opspf:2019} or segment-routing
policies on the same StarryNet topologies would clarify when centralized
install wins over orbit-predictive link-state flooding, and whether hybrid
designs (OSPF data plane with SDN-triggered make-before-break at handover
ticks) capture most of the benefit with less operational coupling.

\paragraph{Traffic engineering and failure models beyond bulk loss.}
Future scenarios could include capacity-aware weighting (not only delay),
time-varying demand hotspots, partial link degradation (elevated loss without
total failure), and targeted ISL failures rather than whole-satellite blackouts.
Each variant stresses a different facet of the graph snapshot model
(Section~\ref{sec:orbitgraph}).

\paragraph{Damage and recovery with SDN-aware ordering.}
Section~\ref{sec:analysis-damage} shows OrbitGraph lagging OSPF after
30\% satellite loss because updates are reactive and bulk.
Promising directions include precomputed backup next-hops for known failure
domains, staged install (restore connectivity to ground first, then refine
intra-orbit paths), and ``soft damage'' detection that triggers recomputation
before \texttt{netem} loss reaches 100\%.

\subsection{Toward real constellation scale}

\paragraph{Scaling beyond $10{\times}10$.}
Container emulation allocates one Linux network stack per node; our workstation
campaign tops out near 102 nodes before memory and CPU become limiting factors,
consistent with observations in the StarryNet literature~\cite{lai2023starrynet}.
Reaching Starlink-class scale (thousands of satellites) requires a tiered
strategy rather than a flat all-container grid:
\begin{itemize}
  \item \textbf{Hierarchical control.} Partition the constellation into orbital
        planes or regions; run OrbitGraph per partition with border-route
        summarization at domain gateways, mirroring terrestrial SDN
        hierarchies~\cite{gregory:satellite-routing:2022}.
  \item \textbf{Multi-machine emulation.} StarryNet supports distributed
        deployment across hosts; sharding containers by plane would raise the
        feasible node count while preserving real protocol stacks on each
        emulated satellite.
  \item \textbf{Hybrid emulation and analytics.} Combine packet-level emulation
        on a representative subgraph (e.g., one plane plus adjacent GS links)
        with flow-level or statistical simulators such as
        Hypatia~\cite{kassing2020exploring} or
        StarPerf~\cite{lai2020starperf} for full-constellation stress tests,
        using the same $G_t$ abstraction in both tiers.
\end{itemize}

\paragraph{Real ephemeris and operator timelines.}
Our grids use Walker-style $O{\times}S$ layouts with Starlink-inspired altitude
and inclination but not live TLE feeds.
Feeding OrbitGraph with public two-line element sets or operator-published
contact schedules would validate proactive handover against real handover
cadences and eclipse intervals, and quantify how far ahead the controller must
commit routes for zero-outage guarantees.

\paragraph{Hardware and over-the-air testbeds.}
Emulation proves semantics; hardware proves timing under RF and processing
constraints.
A natural progression is to port the OrbitGraph controller to a small
software-defined radio or satellite testbed (e.g., a few USRP nodes or an
existing LEO lab network), pushing routes to Linux routers that mirror the
kernel FIB programming path validated here, before any on-orbit trial.

\subsection{Summary}

Closing the install-time gap, broadening workloads and failure models, and
climbing the scaling ladder through hierarchy and hybrid simulation are the
main levers to move from proof-of-concept to operator-relevant deployment.
The core idea---treat the LEO network as a sequence of graphs, compute
shortest paths centrally, and time delta installs against known topology
events---remains unchanged; future work fills in the engineering surround
that our current emulation study deliberately simplified.
